\chapter{Conclusion and Future Work}
\label{chap:Five}

\section{Summary of Achievements}
\label{sec:summary}

This thesis has documented the initial phases of building a "Truth Scanner"\textemdash a neuro-symbolic pipeline for translating informal mathematical definitions into formally verified Lean 4 code. The work completed through Week 4 establishes a solid foundation for automated mathematical formalization:

\begin{itemize}
    \item \textbf{Environment Mastery}: Successfully configured a complete Lean 4 development environment with Mathlib integration and completed the Natural Number Game to gain practical proficiency in tactic-based theorem proving.
    \item \textbf{Systematic Selection}: Identified and taxonomized 10\textendash15 key definitions from Rosen's \textit{Discrete Mathematics}, mapping informal English to formal logical structures.
    \item \textbf{Manual Baseline}: Created a gold standard set of manually formalized definitions that compile successfully in Lean 4, establishing baseline metrics for accuracy and completeness.
\end{itemize}

These accomplishments demonstrate the feasibility of bridging the gap between textbook mathematics and machine-verifiable formal logic.

\section{Contributions}
\label{sec:contributions}

The primary contributions of this work include:

\begin{enumerate}
    \item \textbf{Verified Formalization Baseline}: A collection of manually translated mathematical definitions from a standard undergraduate textbook, serving as ground truth for evaluating automated approaches.
    \item \textbf{Translation Methodology}: A documented systematic process for decomposing informal mathematical language into Lean 4 dependent types.
    \item \textbf{Challenge Analysis}: Identification of common pitfalls and ambiguities encountered when formalizing textbook mathematics, including implicit assumptions and notation mismatches.
\end{enumerate}

\section{Future Work}
\label{sec:future-work}

The next phases of this research will implement and evaluate the automated "Truth Scanner" pipeline:

\subsection{Week 5: Pipeline Development (Phase I)}
Develop a Python script using the OpenAI/Gemini API to automate Lean 4 translation:
\begin{itemize}
    \item Design prompts that take English mathematical text as input and output Lean 4 code snippets
    \item Implement basic error handling and code generation
    \item \textbf{Deliverable}: A script that takes English math text and outputs Lean code snippets
\end{itemize}

\subsection{Week 6\textendash7: The Repair Loop}
Implement the self-correction mechanism that addresses LLM hallucinations:
\begin{itemize}
    \item Feed Lean compiler error messages back into the LLM prompt
    \item Iterate until the code compiles or exceeds maximum retry attempts
    \item Track success rates, iteration counts, and error patterns
    \item \textbf{Deliverable}: An end-to-end pipeline with repair loop functionality and preliminary evaluation metrics
\end{itemize}

\subsection{Week 8: Evaluation and Analysis}
Conduct comprehensive testing of the automated pipeline:
\begin{itemize}
    \item Compare LLM-generated translations against the manual baseline
    \item Analyze success rates, common failure modes, and repair loop effectiveness
    \item Categorize errors to identify systematic weaknesses in LLM mathematical reasoning
    \item \textbf{Deliverable}: A detailed analysis report with visualizations and statistics
\end{itemize}

\subsection{Week 9: Documentation and Finalization}
Complete the thesis documentation:
\begin{itemize}
    \item Finalize all chapters with complete results and discussion
    \item Create a GitHub repository containing all Lean files and the pipeline script
    \item Write a report titled "Where the Textbook was Ambiguous" documenting discovered gaps
    \item \textbf{Deliverable}: A complete thesis and public code repository
\end{itemize}

\section{Long-Term Vision}
\label{sec:vision}

The ultimate goal of this research is to create a "Wikipedia of Truth" where mathematical claims are automatically backed by machine-checked proofs. By debugging educational materials at scale, this work aims to:

\begin{itemize}
    \item Improve the precision and clarity of mathematical textbooks
    \item Detect subtle logical errors that human reviewers might miss
    \item Accelerate the formalization of mathematical knowledge
    \item Evaluate the current capabilities\textemdash and limitations\textemdash of LLMs in maintaining mathematical rigor
\end{itemize}

As Large Language Models continue to advance, the "Truth Scanner" pipeline represents a crucial step toward trustworthy AI-assisted mathematics, where creative language understanding meets uncompromising logical verification.
